\documentclass[11pt]{article}
\usepackage{../homework}

\profname{Dr. David Singer}
\classcode{MATH 465 Differential Geometry}
\assignnum{Assignment 06}
\duedate{26 October 2026}

\begin{document}
\begin{problem}{4.3.2}
  Show that if \( \mathbf{x} \) is an isothermal parametrization, that is, \( E = G = \lambda(u,v) \) and \( F = 0 \), then
  \begin{equation}
  \label{eq:1}
  K = - \frac{1}{2\lambda}\Delta(\log \lambda),
\end{equation}
where \(\Delta\varphi\) denotes the Laplacian of the function \(\varphi\).
Conclude that when \( E = G = (u^2 + v^2 + c)^{-2} \) and \( F = 0 \), then \( K = \text{const.} = 4c \).
\end{problem}

\begin{problem}{4.3.3}
  Verify that the surfaces
  \begin{equation}
  \label{eq:2}
  \begin{aligned}
    \mathbf{x}(u,v) &= (u\cos v, u\sin v, \log u), \quad u > 0 \\
    \mathbf{\bar{x}}(u,v) &= (u\cos v, u \sin v, v),
  \end{aligned}
\end{equation}
have equal Gaussian curvature at the points \( \mathbf{x}(u,v) \) and \( \mathbf{\bar{x}}(u,v) \) but that the mapping \( \mathbf{\bar{x}} \circ \mathbf{x}^{-1} \) is not an isometry.
This shows that the converse of the Gauss theorem is not true.
\end{problem}

\begin{problem}{4.3.4}
  Show that no neighborhood of a point in a sphere may be isometrically mapped into a plane.
\end{problem}

\begin{problem}{4.3.8}
  Compute the Chrisoffel symbols for an open set of the plane, use the Gauss formula to compute \( K \) in both cases.
  \begin{enumerate}[label=(\alph*)]
    \item In Cartesian coordinates.
    \item In polar coordinates.
  \end{enumerate}
\end{problem}

\end{document}
